\documentclass{article}

\usepackage[dblblindworkshop]{neurips_2026}
\workshoptitle{Foundation Models for Temporal Systems (FMTS): From Forecasting to World Modeling}

\usepackage[utf8]{inputenc}
\usepackage[T1]{fontenc}
\usepackage{hyperref}
\usepackage{url}
\usepackage{booktabs}
\usepackage{amsfonts}
\usepackage{amsmath}
\usepackage{nicefrac}
\usepackage{microtype}
\usepackage{xcolor}
\usepackage{graphicx}
\usepackage{siunitx}

\title{Forecasting Is Not World Modeling: Accuracy, Action Fidelity, and an Evaluation Matrix for a Task-Critical Thermal Process}

\author{Anonymous}

\begin{document}
\maketitle

\begin{abstract}
Task-critical industrial processes demand more from a temporal model than accurate
forecasts. A model promoted to the inner loop of a controller must respond to candidate
actions with the right sign, magnitude, and transient shape, and it must know when it
cannot. The question arose on a live unit where a monitoring forecast had been wired
into the control chain: a 180-step-ahead predictor whose 90-second-ahead prediction
corrects the temperature setpoint as differential feedforward. That system forecasts
well, and it barely responds to its own valves. We study this gap on a once-through
boiler in dry operation, using 82 days of dual-side field data. A deep forecaster
attains 0.7--0.9~\si{\degreeCelsius} MAE at 18 steps, yet its valve-step responses span
0.00004--0.040~\si{\degreeCelsius} with unstable direction across folds, valves, and
seeds---functionally disabling the action channel. A physics-structured world model
trades accuracy for an action channel that a first training exposed as directionally
wrong at 60 steps---an inversion we localize to one feedback term---and that a
contracted repair batch then corrected: on the repaired seed the direction holds on
100\% of windows (mean $-0.228$~\si{\degreeCelsius} per $+5\%$ valve step), with the
remaining seeds pending. Both findings are produced by a discrimination matrix: an
identifiability gate on same-type synthetic data, per-layer structural ablations, and
an action-blindness and shuffle-controlled leakage audit---a protocol whose own false
trigger we report and repair in the loop. Finally we locate the physics model's
accuracy cost: it approaches persistence at the terminal stage (0.16 vs.\
0.09~\si{\degreeCelsius} one-step MAE, beating persistence at 18 steps) while the
upstream stage fails 38$\times$ worse than persistence, and the earlier monotone
regime bias collapses to $\pm$0.2~\si{\degreeCelsius}---a training artefact, not a
parameter-count one. Our results support verifiability, not accuracy, as the
qualifying property of a world model, and a two-tier architecture in which a learned
proposer is bounded by a certifiable reduced-order model. Matrix verdicts are
single-side (side A); the side-B matrix is deferred to a later stage.
\end{abstract}

\section{Introduction}

A temporal model that forecasts well is not automatically a world model. World models
must maintain a state, imagine futures under interventions, and support decisions
\cite{ha2018world,hafner2023dreamerv3}. Surveys frame the same tension as
``understanding vs.\ predicting'' \cite{ding2024understanding}; this paper
instantiates it in a fault-critical setting. Closed-loop identification warned
decades ago that passively learned dynamics absorb the confounded structure of the
controller that generated the data \cite{forssell1999closedloop}, and recent
evidence shows per-series evaluation masking the cross-series structure that
downstream decisions require \cite{wan2026forecast}. This paper is a case study of
that tension being encountered in the field, chased to its root, and turned into a
reusable protocol.

\textbf{Where we started.} On a once-through boiler at an industrial plant, a
monitoring forecast was extended toward closed-loop use. The model predicts 180 steps
ahead, but the control system accepts a single connection point; the 90-second-ahead
prediction---the horizon with the best one-step accuracy and the most reliable action
response---is wired into the setpoint-correction loop as differential feedforward,
realizing lead control ahead of the plant's multi-minute response lag. The system
forecasts well. It also barely reacts to its own valves: changing the spray-valve
positions leaves the prediction almost unchanged, and where a response does appear its
direction is unreliable. For a feedforward loop whose physical lever \emph{is} the
valve, this is disqualifying.

\textbf{Why not learn a world model the usual way.} In robot control and autonomous
driving, the standard remedies are interaction-based: online reinforcement learning
explores the action space, offline RL reuses logged interaction or a simulator
\cite{hafner2023dreamerv3}. Neither is available here. Online trial is barred by
safety class: a superheater excursion trips the unit, and exploratory dithering of
spray valves---the discrete probing RL would want---is operationally unacceptable.
Offline training in simulation is barred by a transfer gap the other way: the boiler
is strongly coupled, disturbance-rich, and time-varying, and a simulator faithful
enough that controllers learned inside it transfer to the plant does not exist.
Closed-loop operating records are the only admissible data. The problem this paper
addresses is therefore sharper than generic world-model evaluation: \emph{how to
identify the causal sign and a verified bound on the magnitude of an action channel
from data generated under feedback, when the controller systematically couples the
valve to the very temperature it drives}. This common-cause structure is precisely what the diagnosis
below unpacks.

\textbf{The diagnosis.} Our earlier diagnosis quantified the anomaly in three layers.
First, representation: the forecast did not ignore the valves by architecture---a
model with a zeroed action pathway still reached 0.808~\si{\degreeCelsius} rollout
MAE, and switching from differential to absolute valve inputs enlarged the learned
response by a factor of 32--130. One-step metrics are blind to it by construction: the
plant's true spray-to-outlet lag exceeds 120~s, so a correct one-step response is
nearly zero. Second, timescale: single-step sign regularization therefore not only
failed, it compressed a 120~s physical lag into a spurious 10~s response. Third,
confounding: the controller opens the spray valve when temperature rises, so
``valve-open $\leftrightarrow$ temperature-rising'' is a statistical fact of the data
and a physical falsehood.

\textbf{The illusion.} A constrained gray-box branch appeared to fix the direction:
its action response was sign-consistent across the evaluation. A 1~s-resolution audit
exposed what it had learned. Setpoint changes are rare (0.76\% of steps) and the
exogenous valve signal is rarer still (0.2\% of steps; 99.8\% of valve motion is
feedback-driven). The branch's ``correct'' sign was closed-loop \emph{tracking}: the
setpoint rises, the controller closes the spray valve, the temperature rises---the
model learned to track the setpoint through the controller, not the plant's physical
valve-to-temperature pathway. Direction had been faked once. Every claim made for the
physics line in this paper is therefore gated by an audit, not by its architecture.

\textbf{Two lines, each failing alone.} The pure forecasting line (a high-capacity
direct forecaster with an action channel) predicts held-out trajectories well but its
valve responses span 0.00004--0.040~\si{\degreeCelsius} with direction drifting
across folds, valves, and seeds. The pure physics line (a lumped-enthalpy ODE,
descended from the constrained branch, with a residual closure) responds on the right
scale but instantaneously---the plant's response builds over minutes---and its
accuracy is large and initial-state dependent, because the field data carry no
reliable spray-flow, wall-temperature, or enthalpy ground truth. Neither line alone
qualifies as an internal model for control.

\textbf{What follows.} We measure the resulting tension three ways and report it
through a discrimination matrix: an identifiability gate, structural ablations, and
an action-blindness and leakage audit, run by a frozen protocol with verdicts
computed by paired bootstrap and a 2/3 seed rule. The matrix is not a checklist; it
is a process. Its verdict-repair loop caught its own leakage probe misfiring and
protocolized the correction, and one of its convergence diagnostics demoted a
parameter-map research line whose evidence turned out to be a training artefact. Our
contributions are (i) an empirical demonstration, on dual-side field data, that
forecast accuracy does not transfer to world-model qualification; (ii) the
discrimination matrix and the protocolized verdict-repair loop it produced; and (iii)
a located diagnosis of the physics model's accuracy cost, with quantified floors and
a convergent repair agenda.

\section{Testbed and Two Model Lines}

\subsection{The process}

The unit drives steam through a chain of superheater stages interleaved with two
spray attemperators: low-temperature superheater, first spray, platen superheater,
second spray, final superheater. Five temperatures along this chain are observed at
10~s resolution, together with boundary channels (steam flow, coal command, separator
pressure and temperature, feedwater temperature, outlet pressure, spray flow) and two
valve positions. The valves are cross-paired: valve A governs side-B temperatures and
vice versa. Data span 82 days per side, all in dry operation. Fast temperature
increments are small (0.09--0.45~\si{\degreeCelsius} per step), so the process is
strongly autoregressive; load-driven swings are orders of magnitude larger.

\subsection{Line 1: a deep forecaster}

Line 1 is a direct multi-step forecaster (sequence model over a history window with
dedicated heads per horizon step, trained with a Gaussian likelihood), the audited
descendant of the deployed predictor with an explicit action channel. It is
representative of the pure forecasting line rather than a state of the art: the
horizon heads are teacher-forced, not recursive. We state this plainly because the
comparison that follows is about what the forecasting objective rewards, and the
literature suggests recursive foundation forecasters inherit the same closed-loop
failure \cite{wan2026forecast}. Its 18-step MAE on held-out windows is 0.712 (side A)
and 0.927 (side B).

\subsection{Line 2: a physics-structured world model}

Line 2 embeds the process structure: three lumped-enthalpy segments whose dynamics
obey energy conservation with IAPWS-97 property tables, two spray injections with
transport delays, metal-temperature coupling, and learnable parameters (segment
conductances, evaporation gains, spray efficiencies), in the
universal-differential-equation tradition of embedding structure while learning
parameters \cite{chen2018neuralode, rackauckas2020ude}. It descends directly from the
constrained gray-box branch of Section 1: that branch's apparent direction fidelity
was exposed as controller tracking, and the deeper enthalpy structure is the response
to that exposure. An action-blind closure absorbs unmeasured disturbances from state
and whitelisted boundary channels only. A probabilistic observer estimates the
initial state. The model rolls out recursively, so any counterfactual branch is a
consistent trajectory of the same dynamics. Its 18-step MAE is about
3.9~\si{\degreeCelsius} under the amended training budget (side A, seed-averaged;
see Section 4).

\section{The Tension, Measured Three Ways}

\subsection{Action response: one line is functionally blind}

We probe both models with $\pm$2\% valve steps, holding all else fixed, and compare
against a data-anchored mixing reference (spray-flow regressions on the field data,
$dW/dv = 27.8$ and $70.0$~t/h per full travel for the two valves, combined with an
energy balance at the spray station; the reference spans 0.21--0.59~\si{\degreeCelsius}
for the primary valve and 0.53--1.48~\si{\degreeCelsius} for the secondary valve per
$+2\%$, depending on the steam-flow band). Line 1 responds with mean magnitudes of
0.00004--0.040~\si{\degreeCelsius} across folds, valves, and directions, with the
correct sign in only one of eight fold--valve--direction cells at 3/3
consistency---two to four orders of magnitude below the reference, and functionally
blind (Fig.~\ref{fig:evidence}a).

Line 2's first training exposed a different failure. Its one-step response lands on
the reference scale, but partly by construction: the action enters an instantaneous
mixing equation re-evaluated at every step, so the model responds at step one whether
or not the learned dynamics agree, while the real plant's one-step response is
effectively zero and builds over minutes. At 60 steps the direction broke: only
19--34\% of windows ended negative (the direction rule as frozen required 100\%), with
a mean terminal delta of $+0.20$ to $+0.26$~\si{\degreeCelsius}. An ablation
localized the inversion to one named term: the wall-rewetting feedback. More spray
raises the droplet mass, lowering the dry-out factor, raising the rewetting heat flux
into the steam; the learned rewetting gain over-compensated the direct cooling.
Zeroing the rewetting coefficient restored the correct sign on 100\% of windows (mean
$-0.44$~\si{\degreeCelsius})---a pre-repair diagnosis, not a repair.

The contracted repair batch then changed the model, not the probe: a
spray-to-mixing transport-lag state, a sign-and-magnitude contract on the rewetting
feedback, and spray-sensitivity priors anchored to the data-fit $dW/dv$. The
retrained seed-0 model holds the correct direction on 100\% of windows with a mean
terminal response of $-0.228$~\si{\degreeCelsius} per $+5\%$ step, and its 240-step
(quasi-steady) response is $-0.194$~\si{\degreeCelsius} per $+2\%$. This remains
2.7--7.6$\times$ below the mixing reference. We read the gap as a distance to a
nominal linear upper bound, not as demonstrated distortion of the plant: the
reference assumes zero-delay mixing, linear valve-to-flow mapping, and fixed
$\Delta h/c_p$, while the real valve is equal-percentage and the data contain only
narrow-band micro-adjustments (command changes $>$3\% are absent), so the identifiable
quantity is a workpoint-local gain---consistent with the learned spray efficiency
landing at 0.44$\times$ its prior. Magnitude fidelity remains unverified pending a
position-binned local-gain probe. Seeds 1--2: pending (\texttt{<R1-3SEED>}).

\subsection{Counterfactual coherence}

Beyond steady responses, an internal model must carry physically coherent
counterfactual branches. Line 1's teacher-forced heads cannot propagate a perturbation
through time; Line 2's recursive state does so by construction, and its boundary
dependence is whitelisted so that no action leaks into the disturbance path. The
audit of that channel is itself part of the story, told in Section 4: the leakage
probe that first condemned the model was, on a shuffle control, condemning itself.

\subsection{Accuracy: the cost is real, located, and mostly structural}

Line 2 pays a 3--4$\times$ aggregate accuracy price (H18 MAE
3.9~\si{\degreeCelsius} vs.\ 0.7--0.9 for Line 1), but the price is not uniform
across the steam path. At the terminal stage the retrained model is nearly as good
as persistence one step ahead (0.16 vs.\ 0.09~\si{\degreeCelsius} MAE) and beats it
at 18 steps (1.05 vs.\ 1.18); at the upstream stage it is 38$\times$ worse than
persistence (9.4 vs.\ 0.25~\si{\degreeCelsius}), and the model's learned noise scale
confirms it has written this channel off ($\sigma \approx 8.8$~\si{\degreeCelsius}).
Binning residuals by load shows that the strong monotone regime bias present under
the initial training budget ($+2.1$ to $+5.4$~\si{\degreeCelsius} across load bins at
the terminal stage) collapses to $\pm$0.2~\si{\degreeCelsius} after retraining with
the amended budget: level calibration was a training-sufficiency artefact, not a
parameter-count artefact. What remains is structural: the upstream channel's
residual swings $\pm$8.8~\si{\degreeCelsius} across load bins with no monotone
signature, pointing at the enthalpy--temperature inversion at the first segment
(pressure partition assumption) and the anchor coverage (two spray-outlet
temperatures float at initialization), and the first spray-outlet channel retains
44\% regime-correlated variance, the signature of unresolved spray transients
(Fig.~\ref{fig:evidence}b,c). The autoregressive diagnosis generalizes: the process
creeps (0.09--0.45~\si{\degreeCelsius} per step), so large one-step errors are model
jumps, not noise.

\section{A Discrimination Matrix for World-Model Qualification}

The tension above is only meaningful if the measurements are reproducible. We froze an
evaluation protocol, run end-to-end by a single runner against a versioned data
contract, with verdicts computed by paired bootstrap confidence intervals and a 2/3
seed rule. Table~\ref{tab:verdicts} summarizes the side-A verdicts.

\begin{table}[t]
\centering
\small
\begin{tabular}{lll}
\toprule
Unit & Side-A verdict & Seeds \\
\midrule
D-SYN identifiability & SUPPORTED & 3/3 (NLL $-99$ to $-136$) \\
O1 posterior, learned & MIXED & $-30.1$, $-31.4$, $+13.3$ (\%) \\
O1 posterior, hybrid & REJECTED & $-0.3$, $-1.6$, $+3.7$ (\%) \\
T1 closure (conservative) & SUPPORTED & $+5.7$, $+6.9$, $+2.6$ (\%) \\
T1 closure (steam-only) & REJECTED & $-6.4$, $-2.3$, $-3.7$ (\%) \\
T1 latent block & REJECTED & $-1.6$, $+0.3$, $+3.1$ (\%) \\
B1 boundary forecaster & REJECTED & CRPS gaps $1.7$--$2.0$ \\
J1 joint training & SUPPORTED & $+13.5$, $+16.4$, $+33.3$ (\%) \\
R1 blindness & SUPPORTED & 3/3 \\
R1 leakage (shuffle-controlled) & SUPPORTED & $\Delta = 0.6\%$ \\
R1 direction & v0.2 REJECTED; repair-batch seed 0 tentative PASS & 19--34\% $\to$ 100\% \\
\bottomrule
\end{tabular}
\caption{Side-A discrimination matrix verdicts. Percentages are validation-NLL
improvements (CRPS gaps for B1); confidence intervals per unit as in the committed
summary. R1 direction: windows ending negative under a $+5\%$ secondary-valve step
before (three seeds) and after (seed 0) the repair batch.}
\label{tab:verdicts}
\end{table}

\textbf{D-SYN identifiability gate.} The full stack (observer, transition, closure)
trains on synthetic data generated by the same architecture and must recover the
teacher: 3/3 seeds pass with NLL reductions of 99--136, well above the 30\% threshold.
Failures on real data therefore reflect structure or data, not an unidentifiable task.

\textbf{O1 initial-state posterior.} A learned posterior must beat the steady-state
anchor by at least 5\% with confidence. It is MIXED (two of three seeds degrade
significantly, one improves), and the precision-weighted hybrid is REJECTED (no seed
reaches a significant improvement). Under the amended training budget the steady
anchor itself improved enough to erase the posterior's advantage; the learned
posterior's seed instability remains.

\textbf{T1 transition structure.} Nested arms add closure layers and a latent block
against the previous layer. The conservative closure is SUPPORTED (2/3 seeds improve
val NLL by 5.7\% and 6.9\% with confidence; seed 2 by 2.6\%, CI crossing zero). The
steam-only placement and the latent block are REJECTED against it. Under the amended
budget the latent arm converged to the same valley the best seed found before, but
not consistently enough to earn its parameters.

\textbf{B1 boundary forecast.} The learned boundary forecaster is REJECTED: all three
seeds are significantly worse than persistence on CRPS (gaps of 1.7--2.0, CIs
entirely below zero). Boundary channels are strongly autoregressive; a learned
forecaster needs more than a GRU to earn its complexity here---a negative result we
report as part of the matrix, not around it.

\textbf{J1 training scheme.} Joint training of transition and observer against staged
training: SUPPORTED 3/3 (val NLL improvements of 13.5\%, 16.4\%, 33.3\%, all CIs
above zero).

\textbf{R1 action-blind and leakage audit.} The causal-diagnosis language follows
causal discovery for autocorrelated series \cite{runge2020pcmci}, adapted to a plant
where the controller couples the valve to the temperature it drives. The closure must
not respond to valve channels (runtime blindness check: 3/3 pass) and the boundary whitelist excludes the
spray-flow channel. The leakage probe now runs with a shuffle control: an
action-aware probe must beat its own shuffled-action twin, not just the action-blind
baseline. This rule is itself a product of the matrix. The original probe flagged
23.9\% ``leakage'' on the repaired seed; the shuffle control scored 23.2\%, a
$\Delta$ of 0.6\%---the flag was a probe artifact (an underfit 20-epoch budget where
any extra input dimension improves validation fit), and the repaired seed passes the
gate. The direction probe is reported above: v0.2 failed it (19--34\% of windows
ending negative), the repair batch corrected seed 0 (100\%, $-0.228$~\si{\degreeCelsius}
per $+5\%$), and seeds 1--2 are pending (\texttt{<R1-3SEED>}).

Each verdict is computed by the frozen runner, not by hand, and every number in this
section reproduces from the committed artifacts through the frozen replay script
(available on request to qualified reviewers). The matrix doubles as the artifact we
propose as an honest benchmark for world-model claims \cite{ghosh2026rom}: the
proposal is not the checklist but the verdict-repair loop---one gate correction
(shuffle-controlled leakage) and one convergence diagnostic (the regime-bias
collapse) each changed the protocol or the research agenda, and both are documented
in the audit trail. All verdicts are side A only; the side-B matrix is deferred and
the single-side scope is carried through the Limitations.

\begin{figure}[t]
\centering
\includegraphics[width=0.49\textwidth]{figs/fig1_tension.pdf}
\caption{The accuracy--action-fidelity plane (one-step magnitude, log scale). The deep
forecaster is accurate but its valve response spans 0.00004--0.040~\si{\degreeCelsius}
with unstable direction; the physics world model responds on the reference scale at
one step (10--30$\times$ the plant's near-zero one-step response); its 60-step
terminal response sign-inverted before the repair batch and holds direction on the
repaired seed. Shaded band: data-anchored mixing reference for the secondary valve
per $+2\%$.}
\label{fig:tension}
\end{figure}

\begin{figure}[t]
\centering
\includegraphics[width=\textwidth]{figs/fig2_system_matrix.pdf}
\caption{Left: the steam path with five observed temperatures and two cross-paired
spray valves. Right: the discrimination matrix with side-A verdicts.}
\label{fig:system}
\end{figure}

\begin{figure}[t]
\centering
\includegraphics[width=\textwidth]{figs/fig3_evidence.pdf}
\caption{(a) Real secondary-valve events versus model probes: the plant's response
builds over minutes and ends negative on most windows; the physics line's 60-step
terminal response sign-inverted before the repair batch and is sign-consistent after
it; the deep forecaster stays inside $\pm$0.04~\si{\degreeCelsius}. (b) One-step error
by channel: the physics line is near persistence at the terminal stage and 38$\times$
worse upstream. (c) Residual means by load quintile: the terminal-stage bias is flat
after retraining, the upstream stage swings 5--13~\si{\degreeCelsius}.}
\label{fig:evidence}
\end{figure}

\section{Discussion}

The repaired physics line holds the action direction that the first training lost,
and its terminal-stage accuracy is near persistence; the remaining cost is
concentrated and diagnosable. The upstream channel's one-step error points at the
initialization and inversion layer, not the dynamics: anchoring all five observations
at initialization (two spray-outlet temperatures currently float) and replacing the
fixed pressure-partition assumption with a learned or measured one targets most of
the 9.4~\si{\degreeCelsius} directly (\texttt{<FIX1-BRANCH>}: whether this enters the
paper window is decided after the three-seed verdict). Regime-dependent parameter
maps, the tempting first remedy, lost their primary evidence once the training budget
was corrected---the monotone load bias they were meant to absorb was a training
artefact. This ordering illustrates the value of the matrix: one convergence
diagnostic and one direction probe re-prioritize the entire optimization agenda.

The magnitude gap deserves a precise status. Direction fidelity is now demonstrated
on seed 0; magnitude fidelity is not. The model's quasi-steady response is
2.7--7.6$\times$ below the nominal mixing reference, and we argue the reference
overstates the plant at the operating point (zero-delay linear mixing against an
equal-percentage valve fed by narrow-band excitation). That argument is a hypothesis
until the position-binned local-gain probe lands (\texttt{<CF3>}); until then we
claim direction, not magnitude, for the repaired line. Similarly, the abstract's
promise that the model ``knows when it cannot'' is exactly what the calibration
probe must earn (\texttt{<D1>}): the sentence stays only if the interval coverage
checks out.

For deployment we endorse the convergent architecture the reduced-order-model lineage
articulates \cite{ghosh2026rom, benner2015survey, sha2001simplicity}: a learned
proposer bounded by a certifiable reduced-order model that approves, bounds, or rejects
each proposed action within its certified envelope. Planning with imperfect models is a
documented failure mode \cite{jafferjee2021hallucinating, wang2022doublecheck}; the
two-tier architecture is the constructive answer. The discrimination matrix is the
qualification gate for that proposer, and it is, to our knowledge, the first
qualification protocol aimed at the setting RL cannot serve: causal action-channel
identification from confounded closed-loop records of a safety-class plant, where the
only admissible exploration is the plant's own operation. Forecasting remains
necessary; it is not sufficient.

\textbf{Limitations.} Line 1 is teacher-forced, not a recursive foundation forecaster;
we benchmark no recursive forecaster here, and the external evidence we cite concerns
per-series evaluation hiding cross-series structure \cite{wan2026forecast} rather
than recursive rollout under feedback---the recursive case should be re-measured with
one. Data are from a single
unit in dry operation; two-phase regimes are unobserved. The action-response probe
covers valve steps at fixed load; larger interventions and load-coordinated moves
deserve their own protocol. All matrix verdicts are single-side (A); the side-B
matrix is deferred, and no bilateral claim should be read from these results. The
quantitative action-response claims use the data-anchored mixing reference, which is
identified under closed-loop operation and carries that caveat; magnitude fidelity
remains unverified (\texttt{<CF3>}), as does interval calibration (\texttt{<D1>}).
Data and code are not released publicly; the frozen protocol, replay scripts, and
committed artifacts are available on request to qualified reviewers.

\bibliographystyle{plain}
\bibliography{fmts_refs}

\end{document}
